% ============================================================
% CAPITOLO 10 — Progetto 7: Integrazione Full-Stack
% ============================================================
\chapter{Progetto 7: Integrazione Full-Stack}\index{full-stack}\index{monorepo}

\section*{Cosa Costruirai}
Un'applicazione full-stack completamente integrata:
\begin{itemize}
  \item Backend e frontend che comunicano senza errori
  \item Gestione CORS e proxy configurati correttamente
  \item Error handling unificato end-to-end
  \item Un unico \texttt{\_CONTEXT.md} di progetto che governa entrambi
  \item Il pattern Multi-Agente applicato in pratica
  \item Un sistema di sviluppo riproducibile
\end{itemize}

\textbf{Tempo stimato}: 45--60 minuti\\
\textbf{Prerequisito}: Backend (Cap.~6--8) e Frontend (Cap.~9) funzionanti

% ---------------------------------------------------------
\section{La Struttura del Monorepo}

\begin{praticabox}[title={Pratica --- Organizzare la workspace}]
\begin{lstlisting}[language={}]
notes-fullstack/
+-- _CONTEXT.md              <- Contesto globale
+-- SKILL.md                 <- Skill condiviso
+-- PROGRESS.md              <- Memoria persistente
+-- backend/
|   +-- _CONTEXT.md          <- Contesto backend
|   +-- package.json
|   +-- prisma/
|   +-- src/
+-- frontend/
    +-- _CONTEXT.md          <- Contesto frontend
    +-- SKILL.md             <- Skill React
    +-- package.json
    +-- src/
\end{lstlisting}
\begin{lstlisting}[language=bash]
mkdir notes-fullstack
mv notes-api notes-fullstack/backend
mv notes-frontend notes-fullstack/frontend
cd notes-fullstack && git init
\end{lstlisting}
\end{praticabox}

% ---------------------------------------------------------
\section{Il Contesto Globale}

Il \texttt{\_CONTEXT.md} nella root è il documento più importante --- descrive l'intero sistema.

\begin{lstlisting}[language={},title={\texttt{\_CONTEXT.md} di root}]
# Progetto: Notes Full-Stack Application

## Panoramica
Applicazione web full-stack per la gestione di note
personali con autenticazione OAuth 2.0.

## Architettura
\end{lstlisting}

\begin{figure}[H]
\centering
\includegraphics[width=0.85\textwidth]{10-fullstack-dev-arch}
\caption{Architettura Full-Stack in sviluppo}
\end{figure}

\begin{lstlisting}[language={}]
## Contratto API - Formato
Successo: { "success": true, "data": T }
Errore:   { "success": false,
            "error": { "message": "...", "code": "..." } }

## Codici errore
- 401: Non autenticato -> redirect al login
- 403: Non autorizzato -> mostra errore
- 404: Risorsa non trovata -> pagina 404
- 422: Validazione fallita -> errori nei campi
- 500: Errore server -> messaggio generico
\end{lstlisting}

\begin{notebox}
Questo documento è il contratto tra le due parti dell'applicazione. Quando modifichi il backend, l'IA sa cosa aspettarsi dal frontend e viceversa.
\end{notebox}

% ---------------------------------------------------------
\section{Gestione CORS in Sviluppo e Produzione}\index{CORS}

\textbf{In sviluppo}: il proxy Vite (Capitolo~9) inoltra le richieste.

\textbf{In produzione}: il backend serve direttamente i file statici del frontend.

\begin{praticabox}[title={Pratica --- Backend per la produzione}]
\begin{lstlisting}[language={}]
Modifica il backend per servire i file statici del frontend
in produzione.

In app.js:
- Se NODE_ENV === 'production', servi i file da
  ../frontend/dist
- Per route non-API, servi index.html (React Router)
- CORS solo in sviluppo
\end{lstlisting}
\end{praticabox}

% ---------------------------------------------------------
\section{Error Handling End-to-End}

\begin{figure}[H]
\centering
\includegraphics[width=0.85\textwidth]{09-error-flow}
\caption{Flusso degli errori}
\end{figure}

\begin{praticabox}[title={Pratica --- Error handling completo}]
\begin{lstlisting}[language={}]
Implementa gestione errori end-to-end:

1. Backend: middleware globale che converte tutti gli errori
   nel formato standard. Gestisci PrismaClientKnownRequest
   (P2025 = not found), ValidationError, AuthError.

2. Frontend: Axios interceptor che:
   - 401: cancella auth, redirige al login
   - 422: restituisce errori di validazione al form
   - 500: toast di errore generico
   - Errore rete: banner "Connessione persa"

3. Frontend: ErrorBoundary React per errori di rendering.
\end{lstlisting}
\end{praticabox}

\begin{warnbox}
L'error handler del backend \textsc{non} deve mai esporre dettagli interni in produzione. Messaggi come ``Cannot read property `id' of undefined'' o stack trace sono vulnerabilità di sicurezza (Information Disclosure, OWASP~A01~\cite{owasp:top10}).
\end{warnbox}

% ---------------------------------------------------------
\section{Sessione di Sviluppo con \texttt{PROGRESS.md}}

\begin{praticabox}[title={Pratica --- Crea \texttt{PROGRESS.md}}]
\begin{lstlisting}[language={},title={Esempio di file di progresso}]
# Notes Full-Stack -- Progresso

## Stato Attuale
- [x] Backend REST API con Express.js
- [x] Database PostgreSQL con Prisma
- [x] Autenticazione OAuth 2.0
- [x] Frontend React con dashboard note
- [x] CRUD note completo (frontend <-> backend)
- [x] Error handling end-to-end
- [ ] Testing automatizzato (Capitolo 14)
- [ ] Deploy in produzione (Capitolo 15)

## Decisioni Architetturali
- JWT in cookie httpOnly (non localStorage)
- Proxy Vite in dev, static serving in produzione
- Monorepo con contesti separati
- Context API per stato auth (no Redux)

## Problemi Risolti
- CORS: proxy Vite + cors middleware condizionale
- Refresh token: interceptor Axios riprova la richiesta
- OAuth redirect: callback -> backend -> frontend
\end{lstlisting}
\end{praticabox}

% ---------------------------------------------------------
\section{Il Pattern Multi-Agente in Pratica}

Non si tratta di far comunicare agenti tra loro, ma di \textbf{cambiare il ruolo dell'IA} nelle diverse fasi del lavoro.

\subsection{Il pattern Planner-Generator-Evaluator}

\textbf{1.\ Planner (Pianificatore):}
\begin{lstlisting}[language={}]
Analizza il _CONTEXT.md di root. Voglio aggiungere le
"categorie" alle note. NON scrivere codice. Elenca:
1. Quali file del backend vanno modificati e perche'
2. Quali file del frontend vanno modificati e perche'
3. Le migrazioni del database necessarie
4. L'ordine di implementazione consigliato
\end{lstlisting}

\textbf{2.\ Generator (Generatore):}
\begin{lstlisting}[language={}]
Implementa il punto 1 del piano: la migrazione del database.
Aggiungi il modello Category e la relazione con Note.
\end{lstlisting}

\textbf{3.\ Evaluator (Valutatore):}
\begin{lstlisting}[language={}]
Rivedi il codice che hai appena generato. Verifica:
1. La migrazione e' reversibile?
2. Gli endpoint rispettano il formato standard?
3. Il frontend gestisce il caso "nessuna categoria"?
4. Un utente puo' accedere alle categorie di un altro?
Elenca eventuali problemi trovati.
\end{lstlisting}

\begin{tipbox}
Chiedere ``mostrami un diff'' prima di applicare i cambiamenti è una pratica difensiva efficace.
\end{tipbox}

% ---------------------------------------------------------
\section{Script di Sviluppo}

\begin{praticabox}[title={Pratica --- Script npm nella root}]
\begin{lstlisting}[language=json,title={\texttt{package.json} di root}]
{
  "name": "notes-fullstack",
  "private": true,
  "scripts": {
    "dev": "concurrently \"npm run dev:backend\"
            \"npm run dev:frontend\"",
    "dev:backend": "cd backend && npm run dev",
    "dev:frontend": "cd frontend && npm run dev",
    "build": "cd frontend && npm run build",
    "db:migrate": "cd backend && npx prisma migrate dev",
    "db:studio": "cd backend && npx prisma studio"
  },
  "devDependencies": {
    "concurrently": "^8.0.0"
  }
}
\end{lstlisting}
\begin{lstlisting}[language=bash]
npm install
npm run dev  # Avvia backend e frontend insieme
\end{lstlisting}
\end{praticabox}

% ---------------------------------------------------------
\section{Test End-to-End}

\begin{center}
\begin{tabularx}{\textwidth}{l X}
\toprule
\textbf{Test} & \textbf{Verifica} \\
\midrule
Cold start & \texttt{npm run dev} $\rightarrow$ tutto parte senza errori \\
Login Google & Click $\rightarrow$ OAuth $\rightarrow$ Dashboard \\
CRUD completo & Crea $\rightarrow$ Leggi $\rightarrow$ Modifica $\rightarrow$ Elimina \\
Logout & Logout $\rightarrow$ Prova /dashboard $\rightarrow$ Redirect \\
Token scaduto & L'app rinnova automaticamente \\
Backend down & Frontend mostra ``Connessione persa'' \\
Mobile & Layout responsive su emulatore \\
\bottomrule
\end{tabularx}
\end{center}

\begin{checkpointbox}
Se tutti i test passano, l'applicazione full-stack è completa.
\end{checkpointbox}

% ---------------------------------------------------------
\section*{Cosa Hai Costruito Finora}

Fermati un momento. Partendo da zero, in 10 capitoli hai costruito:

\begin{enumerate}
  \item Un \textbf{programma interattivo} Python (Cap.~4)
  \item Un \textbf{task manager CLI} con persistenza (Cap.~5)
  \item Una \textbf{REST API} con Express.js e Swagger (Cap.~6)
  \item Un \textbf{database PostgreSQL} con Prisma ORM (Cap.~7)
  \item Un sistema di \textbf{autenticazione OAuth~2.0} (Cap.~8)
  \item Un \textbf{frontend React} con dashboard e CRUD (Cap.~9)
  \item Un'\textbf{applicazione full-stack integrata} (Cap.~10)
\end{enumerate}

Tutto \textbf{senza scrivere una riga di codice manualmente}.

% ---------------------------------------------------------
\section*{Riepilogo}

\begin{center}
\begin{tabularx}{\textwidth}{l X}
\toprule
\textbf{Aspetto} & \textbf{Dettaglio} \\
\midrule
Architettura & Monorepo con contesti separati \\
Contesto globale & \texttt{\_CONTEXT.md} di root con contratto API \\
CORS & Proxy Vite (dev) + static serving (prod) \\
Error handling & End-to-end dal database al toast UI \\
Multi-Agente & Pattern Planner-Generator-Evaluator \\
Memoria & \texttt{PROGRESS.md} per sessioni continuative \\
Script & \texttt{npm run dev} avvia tutto con concurrently \\
\bottomrule
\end{tabularx}
\end{center}
